\begin{definition}[The program of one process]
  \label{def:aba-procnet}
  \lean{PLTS.ABA.GBCA.StageRec}\lean{PLTS.ABA.CoreRec}\lean{PLTS.ABA.Net.StageSideRec}\lean{PLTS.ABA.Net.ProcRec}\lean{PLTS.ABA.Net.NetEvt}\lean{PLTS.ABA.Net.NLab}\lean{PLTS.ABA.Net.netEvtLabels}\lean{PLTS.ABA.Net.AbdyStageStep}\lean{PLTS.ABA.Net.ABAProcStepN}\lean{PLTS.ABA.Net.ABAProcN}\leanok
  \uses{def:aba-labels, def:aba-gbca-vocab, def:aba-core-vocab, def:aba-flat-reading, def:relabel, def:syncproduct}
  Process $j$ runs the flat reading of Definition~\ref{def:aba-flat-reading} at ABDY22's implementation of Definition~\ref{def:aba-gbca-vocab}: the stage message type is $\mathrm{GBCA.Msg}$, the stage record is $\mathrm{GBCA.StageRec}(n)$, and the stage-side rows are the fourteen of \leandecl{PLTS.ABA.Net.AbdyStageStep} --- the eight multicasts of those levels, the delivery that files a received message, the call and its loop, and the three returns.
  The twenty-nine rows around them are the shared ones.
  The state of process $j$ is
  \[
    \mathrm{ProcRec}(n) \;=\;
      \mathrm{CoreRec}(n) \times \mathrm{StageSideRec}(n),
  \]
  its round-loop record --- the control record of Definition~\ref{def:aba-core-vocab} together with the DECIDED receipt rows $decIn(k)$ --- beside its stage-side record.
  The stage-side record holds a finite map $stages$ from rounds to stage records of Definition~\ref{def:aba-gbca-vocab}, each with its received set rows $received set(k)$, and a flag $terminated$ (D22).
  The stage record of round $r$ is read $\mathrm{stage}\,r$: the retained record where $j$ has touched round $r$, the initial stage record elsewhere.
  The map is finite by its type, so $j$ holds finitely many stage records at every point of a run, and the round advance resets none of them.
  Every stage-side row is accordingly round-indexed rather than round-guarded: it reads and writes the stage record of the round its own label tags, whichever round the round loop is in, under the guards of that stage record and $terminated = \mathrm{false}$.
  Process $j$ therefore answers the messages of a round its round loop has passed, and a stage delivery of any round is filed.
  The row $terminate$ fires on $j$'s own return and $2f+1$ DECIDED receipts, and its effect is to set $terminated$ and nothing else, so the retained stage records stand still from that point.
  The DECIDED rows carry no termination guard, so a process that has terminated keeps relaying the payloads it holds.
  The return $retABA$ and the DECIDED relay are guarded by the round-loop record holding an input, so a process that was never called neither returns nor relays (D8).
  The row $gcallLoop$ reads the stage record of round $r$ without writing it, requiring that the record holds an input, and carries no termination guard.
  A process that has terminated at phase $toCallG$ over an uncalled stage record therefore has a row on neither call label and takes no further round-loop step, an excluded region this reading accepts.
  There is no outbox and no copy of the corrupted set or its budget, so every guard of the automaton $\mathrm{ABAProcN}(P,j)$ reads $j$'s own record.
  A corruption replaces $j$'s program (D23).
  The round-loop record carries the flag \leandecl{PLTS.ABA.CoreRec.corrupted}, written by $j$'s own half of the corruption broadcast, \leandeclnamed{failSelf}{PLTS.ABA.Net.FlatProcStep.failSelf}; every other row above is enabled only while that flag is down, so the record stands still from the corruption and $terminate$ in particular never fires afterwards.
  In place of those rows the replaced program is the single self-loop \leandeclnamed{corruptedIdle}{PLTS.ABA.Net.FlatProcStep.corruptedIdle}, taken on every label other than $\tau$ and the labels of $\mathrm{actsAt}\,j$ --- the labels on which $j$ would act on its own sub-protocol messages --- where it has no row at all.
  Those messages enter through the handshake rows of D11 and the injections $\mathrm{byzG}$ and $\mathrm{byzD}$.
  The visible return is the exception: the network's own $\mathrm{retByz}$ row requires $j \in F$ and no DECIDED evidence, and the self-loop is its other half.
  It runs over the extended alphabet $\mathrm{NLab}(n) = \mathrm{Lab}(n) \oplus \mathrm{NetEvt}(n)$, where $\mathrm{NetEvt}$ names the rendezvous the shared alphabet cannot: the stage multicast $\mathrm{gsnd}(r,j,m)$ and the stage delivery $\mathrm{gdlv}(r,i,j,m)$; the DECIDED relay $\mathrm{dsnd}(j,b)$ and the DECIDED delivery $\mathrm{ddlv}(i,j,b)$; the fused coin return $\mathrm{retWPub}(r,id,c,b)$, which carries beside the coin the payload to be published (D10); the graded-agreement call $\mathrm{gcallLoop}(r,id,b)$ against an already-opened stage record; and the five Byzantine handshake rows $\mathrm{byzCallG}$, $\mathrm{byzCallGLoop}$, $\mathrm{byzRetG}$, $\mathrm{byzCallW}$ and $\mathrm{byzRetW}$ (D11).
  The set $\mathrm{netEvtLabels}(n)$ of these labels is what the composition hides, the silent label of $\mathrm{NLab}(n)$ being that of $\mathrm{Lab}(n)$.
  Every transition is Dirac, and $terminate$ is the one row of the table that fires on $\tau$.
  While $j$'s program stands unreplaced, every label of either alphabet has a row of process $j$, the participant's or an idle self-loop, with two exceptions: $\mathrm{byzCallG}(r,k,b)$ and $\mathrm{byzRetG}(r,k,out,bnd)$ have no row at the process $k$ they name (D11, D22), and a corrupted process's stage messages enter through the network's own $\mathrm{byzG}$ injection instead.
  Once the program is replaced the labels of $\mathrm{actsAt}\,j$ join those two (D23).
  Full synchronisation therefore lets a label's owner move while every other program stands.
  Full rule table, forty-three rows --- twenty-nine shared, fourteen of ABDY22's: \leandecl{PLTS.ABA.Net.ABAProcN}.
\end{definition}
